% =====================================================================
%  Preámbulo del expediente técnico — Centrales Termoeléctricas
%  Trabajo 01: Equipamiento de una central a vapor
%  Compilación prevista: pdfLaTeX + BibTeX (natbib)
% =====================================================================

% ---------- Codificación e idioma ----------
\usepackage[utf8]{inputenc}
\usepackage[T1]{fontenc}
\usepackage[spanish,es-tabla,es-nodecimaldot]{babel}
\usepackage{lmodern}
\usepackage{microtype}

% ---------- Geometría y tipografía ----------
\usepackage[a4paper,top=3cm,bottom=2.5cm,left=3cm,right=2.5cm,
            headheight=15pt]{geometry}
\usepackage{setspace}
\onehalfspacing
\setlength{\parindent}{1.25em}
\setlength{\parskip}{0.35em}

% ---------- Matemáticas y unidades ----------
\usepackage{amsmath}
\usepackage{amssymb}
\usepackage{siunitx}
\sisetup{detect-all,per-mode=symbol,range-units=single}

% ---------- Tablas y figuras ----------
\usepackage{graphicx}
\usepackage{xcolor}
\usepackage{booktabs}
\usepackage{array}
\usepackage{tabularx}
\usepackage{longtable}
\usepackage{multirow}
\usepackage{float}
\usepackage[font=small,labelfont=bf,justification=centering]{caption}
\usepackage{subcaption}

% ---------- Listas ----------
\usepackage{enumitem}
\setlist[itemize]{topsep=2pt,itemsep=1pt,parsep=0pt}
\setlist[enumerate]{topsep=2pt,itemsep=1pt,parsep=0pt}

% ---------- Encabezados y pies ----------
\usepackage{fancyhdr}
\pagestyle{fancy}
\fancyhf{}
\fancyhead[L]{\small\nouppercase{\leftmark}}
\fancyhead[R]{\small\thepage}
\renewcommand{\headrulewidth}{0.4pt}
\fancypagestyle{plain}{%
  \fancyhf{}%
  \fancyfoot[C]{\small\thepage}%
  \renewcommand{\headrulewidth}{0pt}%
}

% ---------- Títulos de capítulo ----------
\usepackage{titlesec}
\titleformat{\chapter}[display]
  {\normalfont\bfseries\Large}
  {\filleft\MakeUppercase{\chaptertitlename}~\thechapter}
  {1ex}
  {\titlerule\vspace{1ex}\filright}
  [\vspace{1ex}\titlerule]
\titlespacing*{\chapter}{0pt}{-10pt}{30pt}
\titleformat{\section}{\normalfont\large\bfseries}{\thesection}{0.8em}{}
\titleformat{\subsection}{\normalfont\normalsize\bfseries}{\thesubsection}{0.8em}{}

% ---------- Cajas de estado y notas ----------
\usepackage[most]{tcolorbox}
\newtcolorbox{cajaestado}[1][]{%
  enhanced,breakable,
  colback=orange!4!white,
  colframe=orange!55!black,
  fonttitle=\bfseries\small,
  title={Estado del capítulo},
  left=6pt,right=6pt,top=4pt,bottom=4pt,
  boxrule=0.6pt,
  #1
}
\newtcolorbox{cajanota}[1][]{%
  enhanced,breakable,
  colback=blue!3!white,
  colframe=blue!45!black,
  fonttitle=\bfseries\small,
  left=6pt,right=6pt,top=4pt,bottom=4pt,
  boxrule=0.5pt,
  #1
}

% ---------- Bibliografía (autor-año) ----------
\usepackage[round,authoryear]{natbib}
\bibliographystyle{plainnat}
\setcitestyle{authoryear,round,semicolon,aysep={,},yysep={;}}

% ---------- Referencias cruzadas ----------
\usepackage[
  colorlinks=true,
  linkcolor=black,
  citecolor=blue!45!black,
  urlcolor=blue!45!black,
  pdftitle={Equipamiento de una central a vapor},
  pdfauthor={Mamani Galindo Renzo Gabriel},
  pdfsubject={Centrales Termoelectricas - Expediente tecnico},
  pdfkeywords={central a vapor, ciclo Rankine, caldera, turbina, generador}
]{hyperref}
\usepackage{url}

% ---------- Datos del documento ----------
\newcommand{\titulotrabajo}{Equipamiento de una central a vapor}
\newcommand{\subtitulotrabajo}{Expediente técnico de ingeniería}
\newcommand{\autorproyecto}{Mamani Galindo Renzo Gabriel}
\newcommand{\universidad}{Universidad Nacional del Altiplano}
\newcommand{\escuelaprofesional}{Escuela Profesional de Ingeniería Mecánica Eléctrica}
\newcommand{\asignatura}{Centrales Termoeléctricas}
\newcommand{\docente}{Paredes Pareja Walter Oswaldo}
\newcommand{\seccion}{Grupo A}
\newcommand{\grupo}{Grupo 7}
\newcommand{\ciudad}{Puno, Perú}
\newcommand{\periodo}{2026}

% ---------- Comandos de trazabilidad ----------
% Uso: \srcref{SRC-001} para referir una fuente registrada.
\newcommand{\srcref}[1]{\textsf{[#1]}}
% Uso: \hipotesis{...} para marcar un supuesto de cálculo.
\newcommand{\hipotesis}[1]{\textcolor{orange!70!black}{\textbf{[HIPÓTESIS:} #1\textbf{]}}}
% Uso: \ejemplo{...} para marcar un ejemplo académico.
\newcommand{\ejemplo}[1]{\textcolor{blue!60!black}{\textbf{[EJEMPLO ACADÉMICO:} #1\textbf{]}}}

% ---------- Profundidad de numeración ----------
\setcounter{secnumdepth}{3}
\setcounter{tocdepth}{2}
